\chapter{从物种方程到可以检验的加热账本}
\section{未知量与矩的定义}
取单电荷离子$q_i=+e$、电子$q_e=-e$，$e>0$。先不施加准中性。密度$n_i,n_e$、平均速度$\vct u_i,\vct u_e$和中心压力张量$\vct P_i,\vct P_e$分别定义为
\begin{align}
n_i&=\int f_i\dd^3v,&n_e&=\int f_e\dd^3v,\\
\vct u_i&=n_i^{-1}\int\vct v f_i\dd^3v,&\vct u_e&=n_e^{-1}\int\vct v f_e\dd^3v,\\
P_{i,ab}&=m_i\int c_{i,a}c_{i,b}f_i\dd^3v,&P_{e,ab}&=m_e\int c_{e,a}c_{e,b}f_e\dd^3v,
\end{align}
其中$\vct c_i=\vct v-\vct u_i$、$\vct c_e=\vct v-\vct u_e$。下标$a,b,c$表示笛卡尔分量，重复指标求和。三阶中心矩和热流的关系为
\begin{equation}
Q_{s,abc}=m_s\int c_{s,a}c_{s,b}c_{s,c}f_s\dd^3v,
\qquad q_{s,a}=\frac12 Q_{s,abb},\quad s=i,e.
\end{equation}
这里$s$只在已分别定义两物种之后缩写相同结构。压力单位Pa，$Q$与热流均为$\mathrm{W\,m^{-2}}$，但一个是三阶张量，一个是向量，不能互换。文献对平行三阶矩与平行能流可能使用不同二分之一约定，本报告以本式为准。

定义$p_s=\operatorname{tr}\vct P_s/3$、$T_s=p_s/n_s$，温度以能量单位表示；$T_s/k_B$才是K。内能$U_s=\operatorname{tr}\vct P_s/2$，流能$K_s=m_sn_su_s^2/2$。这些定义不要求分布为Maxwell型，因此非热多束也能贡献二阶矩。\textbf{有限矩模型并不唯一确定热核心、尾部及精细分布熵。}

\section{分别列出的完整非相对论两物种方程}
无反应、无辐射及无外源时，连续性为
\begin{align}
\partial_t n_i+\nabla\cdot(n_i\vct u_i)&=0,\\
\partial_t n_e+\nabla\cdot(n_e\vct u_e)&=0.
\end{align}
动量方程为
\begin{align}
m_in_iD_i\vct u_i&=en_i(\E+\vct u_i\times\B)-\nabla\cdot\vct P_i+\vct R_i,\\
m_en_eD_e\vct u_e&=-en_e(\E+\vct u_e\times\B)-\nabla\cdot\vct P_e+\vct R_e,
\end{align}
其中$D_i=\partial_t+\vct u_i\cdot\nabla$，$D_e=\partial_t+\vct u_e\cdot\nabla$；$\vct R_s$是碰撞动量交换密度。两套内能方程是
\begin{align}
\partial_tU_i+\nabla\cdot(U_i\vct u_i+\vct q_i)&=-\vct P_i:\nabla\vct u_i+H_i,\label{eq:Ui}\\
\partial_tU_e+\nabla\cdot(U_e\vct u_e+\vct q_e)&=-\vct P_e:\nabla\vct u_e+H_e.\label{eq:Ue}
\end{align}
此处$H_s$是碰撞或明确给定的内能源密度，单位$\mathrm{W\,m^{-3}}$。两个方程中的速度不能擅自统一；电子压缩、对流和焓流通常沿电子流。

全电磁方程及约束为
\begin{align}
\partial_t\B&=-\nabla\times\E,&\epsilon_0\partial_t\E&=\mu_0^{-1}\nabla\times\B-\J,\\
\nabla\cdot\B&=0,&\epsilon_0\nabla\cdot\E&=e(n_i-n_e),\\
\J&=e(n_i\vct u_i-n_e\vct u_e).
\end{align}
五矩取每物种$n_s$、三个动量和一个总粒子能量$K_s+U_s$，假定$\vct P_s=p_s\vct I$并指定热流、碰撞；两物种共有十个流体变量，另有六个电磁演化分量及散度约束。十矩以六个对称压力分量替代标量能量，每物种共十个矩；内能是压力迹，不能再额外独立推进一个不协调的标量能量方程。\citep{Wang2015,TenBarge2019MMS}

\section{Hall约化删去什么、保留什么}
在$\lambda_D/L\ll1$、低频弱电荷分离条件下令$n_i\simeq n_e=n$，在电磁传播时间远短于演化时间的排序下忽略位移电流。于是$\J=\nabla\times\B/\mu_0$且
\begin{equation}\vct u_e=\vct u_i-\frac{\J}{en}.\end{equation}
把电子动量方程除以$-en$再代入，得
\begin{equation}
\E+\vct u_i\times\B=
\frac{\J\times\B}{en}-\frac{\nabla\cdot\vct P_e}{en}
-\frac{m_e}{e}D_e\vct u_e+\frac{\vct R_e}{en}.
\label{eq:ohm}
\end{equation}
令$m_eD_e\vct u_e$在目标尺度上可忽略，就得到无电子惯性的Hall形式；若再取标量压力，则失去非旋转对称压力支撑。电子惯性的删除必须以频率和层厚排序说明，不能仅凭总体质量比小就在电子层删除。准中性也不等于$\vct u_i=\vct u_e$。

质心速度$\vct u=(m_i\vct u_i+m_e\vct u_e)/(m_i+m_e)$下，两物种总流能严格分为
\begin{equation}
K_i+K_e=\frac{\rho u^2}{2}+\frac{m_im_e}{2e^2(m_i+m_e)n}J^2,
\quad\rho=(m_i+m_e)n.
\end{equation}
后一项是相对漂移能。若保留电子惯性却漏记相应能量与通量，可能把可逆漂移储能变化误归为加热。碰撞少时电子惯性提供交换和储存通道，本身不是正定熵源。

\section{电功怎样经过流能进入内能}
利用连续性将动量方程点乘本物种速度，再用
$-\vct u\cdot\nabla\cdot\vct P=-\nabla\cdot(\vct P\cdot\vct u)+\vct P:\nabla\vct u$，分别得到
\begin{align}
\partial_tK_i+\nabla\cdot(K_i\vct u_i+\vct P_i\cdot\vct u_i)
 &=en_i\vct u_i\cdot\E+\vct P_i:\nabla\vct u_i+\vct u_i\cdot\vct R_i,\\
\partial_tK_e+\nabla\cdot(K_e\vct u_e+\vct P_e\cdot\vct u_e)
 &=-en_e\vct u_e\cdot\E+\vct P_e:\nabla\vct u_e+\vct u_e\cdot\vct R_e.
\end{align}
因此压力功在流能和内能之间交换；电场功先进入物种总粒子能，而非直接全部进入温度。与Poynting定理相加后，电功也恰好抵消。封闭碰撞交换必须满足
\begin{equation}
\vct R_i+\vct R_e=0,\qquad
H_i+H_e+\vct u_i\cdot\vct R_i+\vct u_e\cdot\vct R_e=0.
\end{equation}
外部驱动、辐射、电离或开放边界另记，不能靠调整两物种热源使残差消失。

在准中性分解中$\E_H=\J\times\B/(en)$，所以$\J\cdot\E_H=0$。但$en\vct u_i\cdot\E_H$和$-en\vct u_e\cdot\E_H$可分别非零且互相抵消。这一恒等式只涉及欧姆律中特定分量，不允许把实测面内电场全部指定为Hall项。Hall场能重排物种能量路径，却不自动指定最终不可逆热化。

\section{七种量及边界通量}
\begin{longtable}{P{27mm}P{66mm}P{55mm}}
\toprule 名称 & 操作定义 & 易混淆之处\\\midrule\endfirsthead
\toprule 名称 & 操作定义 & 易混淆之处\\\midrule\endhead
物种电功 & $W_s=\int\dd t\int_Vq_sn_s\vct u_s\cdot\E\dd V$ & 含有序流储存和输运，不等于热。\\
流能 & $\int_VK_s\dd V$ & 需指定实验室系及物种平均流。\\
内能 & $\mathcal U_s=\int_V\operatorname{tr}\vct P_s/2\dd V$ & 含未分辨多束的二阶矩。\\
温升 & $\Delta T_s$或$\Delta T_{s,ab}$，指定位置和时间 & 密度和粒子交换改变温度与能量的对应。\\
焓流及压力输运 & 总粒子能流中的$U_s\vct u_s+\vct P_s\cdot\vct u_s$ & 标量时为$(U_s+p_s)\vct u_s$；各向异性时不能只用$5p_s\vct u_s/2$。\\
非热成分 & 指定核心拟合、阈值及参考分布后定义 & 五矩或十矩本身不唯一恢复尾部能量。\\
不可逆热化 & 指定碰撞、黏性、粗粒化或矩熵的定义及收敛 & 正压力功不保证正熵产，温升不保证热化。\\
\bottomrule\end{longtable}

由式\eqref{eq:Ui}积分得到
\begin{equation}
\Delta\mathcal U_i=\int\dd t\int_V(-\vct P_i:\nabla\vct u_i+H_i)\dd V
-\int\dd t\oint_{\partial V}(U_i\vct u_i+\vct q_i)\cdot\dd\vct A.
\end{equation}
内能账本中的对流能流和总能账本中的焓流是两种重排，不可混用后重复加$p\vct u$。$W_i/W_e$、$\Delta\mathcal U_i/\Delta\mathcal U_e$、$T_i/T_e$和所选不可逆热量比应分别输出；分母接近零或变号时应直接报告各项而非稳定“份额”。

\section{十矩新增的方程与热流闭合}
以下中心压力方程对离子、电子各有六个独立分量，分别取$(q_s,m_s,\vct u_s)=(e,m_i,\vct u_i)$和$(-e,m_e,\vct u_e)$：
\begin{align}
D_sP_{s,ab}+P_{s,ab}\partial_cu_{s,c}
+P_{s,ac}\partial_cu_{s,b}+P_{s,bc}\partial_cu_{s,a}
+\partial_cQ_{s,abc}\notag\\
=\frac{q_s}{m_s}\left(\epsilon_{acd}P_{s,cb}+\epsilon_{bcd}P_{s,ca}\right)B_d+C_{s,ab}.
\end{align}
磁项旋转压力张量但迹为零；取半迹返回内能式，其中$H_s=\operatorname{tr}\vct C_s/2$。新增的非对角压力参与欧姆律和应变功，新增的平行、垂直分量参与各向异性。但三阶$Q$仍未确定，十矩并非“无闭合”。

局地松弛常以$\partial_cQ_{s,abc}\simeq\nu_s(P_{s,ab}-p_s\delta_{ab})$建模，$\nu_s\sim v_{t,s}|k_{0s}|$。其迹为零，直接作用是各向同性化；即使矩熵增加，\textbf{该子步并不增加总内能}。热流在真实空间中搬运能量的作用不能仅由这样的零迹项代替。梯度闭合则可写成代表性形式
\begin{equation}
Q_{s,abc}=-\kappa_s n_s\,\mathcal S_{abc}(\partial_aT_{s,bc}),
\quad \kappa_s\sim\chi v_{t,s}/|k_{0s}|,\quad T_{s,ab}=P_{s,ab}/n_s.
\end{equation}
$\mathcal S$的平均或求和规范须按具体代码定义。常系数时$-\nabla\cdot Q$形成扩散，Fourier衰减随$k^2/|k_0|$，有别于局地松弛的固定衰减率。L256扩散的是压力偏离各向同性部分；Ng2020及后继温度梯度形式并不相同，不能只因都叫gradient closure而互换。

非局地Hammett--Perkins型算子近似$|k|v_t$相混响应，通常需Fourier或沿场非局地运算。适用的线性、磁化及几何排序不同；在强导场中用无磁化各向同性扩散可能过度跨场输运。闭合变得更复杂，并不自动保证压力正定、非线性稳定或真实碰撞熵产。\citep{Wang2015,Ng2017,L256,Ng2020,Closure2026}

\section{各向异性与多流体怎样改变加热解释}
旋转对称压力为$\vct P=p_\perp\vct I+(p_\parallel-p_\perp)\vct b\vct b$，内能$U=p_\perp+p_\parallel/2$，故
\begin{equation}
-\vct P:\nabla\vct u=-p_\perp\nabla\cdot\vct u-(p_\parallel-p_\perp)\vct b\vct b:\nabla\vct u.
\end{equation}
相同压缩率下，沿场应变使各向异性模型与标量模型有不同能量交换。无热流、磁化且满足双绝热排序时$D(p_\parallel B^2/n^3)=0$、$D(p_\perp/(nB))=0$，从而$T_\parallel\propto n^2/B^2$、$T_\perp\propto B$。这是可逆变形极限，不能把两个温度的分离直接命名为不可逆加热。俘获电子状态方程在此极限与通过粒子的等温响应之间连接，边界热库是必要信息。\citep{Wetherton2019,Wetherton2021}

再考虑同一物种分成群$a$，令$\vct w_a=\vct u_a-\vct u$。由展开中心矩得到
\begin{align}
\vct P&=\sum_a(\vct P_a+\rho_a\vct w_a\vct w_a),\\
Q_{ijk}&=\sum_a[Q_{a,ijk}+w_{a,i}P_{a,jk}+w_{a,j}P_{a,ik}+w_{a,k}P_{a,ij}
+\rho_aw_{a,i}w_{a,j}w_{a,k}].
\end{align}
即使每群内部热流很小，总三阶矩也可来自群间漂移。F635用这一思路给出电子多流体闭合；它增加的是电子群标签与相对流动，并非新增自洽离子加热。一个总电子流体在数学上仍可表示这些低阶矩，问题在于现有闭合是否能够从演化变量可靠预测它们。\citep{F635}
